\scanpage{44}
\mspage{1}
\uline{Teichmüller groupoids}.

\textcircled{1} Topological case.

The simplest variant is
\[
\mathcal T^+_{\mathrm{top}}\quad\text{or}\quad\mathcal T_{\mathrm{top}}
\]
(\uuline{oriented} \uline{topological} \uline{Teichmüller groupoid})

\uline{objects} : compact oriented surfaces \uline{with boundary} $(\Sigma)$\unsure{,|.}

\uline{arrows} : isotopy classes of isomorphisms (preserving the orientation)

\margin{In the end, as long as one is not working from the
arithmetic point of view, it seems that it is the big $\mathcal T^\pm$ that will be
used the most; \deleted{\unsure{all|whence}} it is therefore this one that
deserves here the simplest notation $\mathcal T$,
and the oriented \unsure{gr.\ or.} to be called $\mathcal T^+$.}

Immediate variant
\[
\mathcal T_{\mathrm{top}}^\pm
\]
(\uuline{bioriented} \uline{topological} \uline{Teichmüller groupoid}) :
same objects, arrows are the isotopy classes of homeomorphisms
(not necessaril\textsuperscript{y} preserving the orientation).

Let $[\pm1]$ be the groupoid having a single object $e$ (for example $e=\varnothing$)
and with automorphism group $\pm1$. We have a sign morphism
\diagram{\begin{tikzcd}[ampersand replacement=\&]
\mathcal T^\pm_{\mathrm{top}}\arrow[r,"\chi"]\&{[\pm1]}
\end{tikzcd}}
whose \uline{fiber} \deleted{\unsure{is}} (or kernel) is
$\mathcal T_{\mathrm{top}}^+=\mathcal T_{\mathrm{top}}$.
\margin{\figure{\begin{tikzpicture}[x=0.27cm,y=0.30cm,every node/.style={anchor=west,inner sep=0pt,outer sep=0pt,font=\tiny}]
\node at (1.2,0) {\unsure{Cannulated|Caution},};
\node at (0,-1.05) {if one does not restrict};
\node at (-0.1,-2.12) {oneself to connected};
\node at (-0.1,-3.18) {surfaces};
\node at (7.8,-2.72) {!};
\draw (7.18,-1.46) -- (7.12,-3.02)
  .. controls (7.13,-3.47) and (6.28,-3.30) .. (5.64,-3.08);
\end{tikzpicture}}}

The most complex variant is
\[
S_*\mathcal TD_{\mathrm{top}}^+\quad\text{or}\quad
S_*\mathcal TD_{\mathrm{top}},
\]
``\inserted{\uline{oriented}} \uline{special decomposed topological}
\uline{Teichmüller groupoid}''
\margin{43}
